\section{Introduction}\label{sec:intro}

Histological grade is one of the strongest prognostic factors in invasive breast cancer. The
Nottingham (Elston--Ellis) system scores tubule formation, nuclear pleomorphism and mitotic count,
and the resulting grade of 1, 2 or 3 predicts survival independently of tumour size and nodal
status~\citep{elston1991grade,rakha2008nhg,rakha2010grade}. Grade has two practical weaknesses.
First, it is only available after a pathologist has examined tissue. Second, it is only moderately
reproducible: agreement between pathologists typically reaches a $\kappa$ of 0.5--0.7
\citep{boiesen2000grading,longacre2006interobserver}, grading rates differ between laboratories in
national data~\citep{vandooijeweert2020variation}, and the large intermediate group, grade~2, is
prognostically mixed~\citep{sotiriou2006ggi,vandooijeweert2022grading}.

Every patient record, however, already contains variables that are linked to grade: age, menopausal
status, oestrogen receptor (ER), progesterone receptor (PR) and HER2 status, histological type,
tumour size and nodal status. If these variables predict grade reliably, a simple model could fill
gaps in registry and trial data, flag records whose assigned grade is unlikely given the rest of the
record, and provide the baseline that image-based grading tools should beat. Recent work has moved
straight to images, predicting grade from whole-slide images with deep learning and pathology
foundation models~\citep{wang2022deepgrade,kim2025grade,chen2024uni} and from MRI
radiomics~\citep{wang2022duke,wang2026dcemrigrade}. We found no study that measured, on open data
from several countries, how much grade the routine record already carries.

\paragraph{Problem statement}
Three gaps prevent a reliable answer.
\begin{enumerate}[leftmargin=*,itemsep=1pt,topsep=2pt]
\item \emph{No transparent baseline.} Published grade-prediction studies are mostly single-centre,
rarely validate their models externally, and seldom report calibration, which is the property that
most often fails when a model moves to a new hospital~\citep{vancalster2019calibration,riley2024evalpart2}.
It is also unclear whether flexible machine-learning methods add anything over regression, which
systematic reviews suggest they usually do not~\citep{christodoulou2019lr}.
\item \emph{Hidden leakage.} Public breast cancer datasets include variables that contain grade by
construction, such as the Nottingham Prognostic Index and the three Nottingham component scores.
Using them as predictors, or fitting preprocessing steps before splitting the data, lets information
about the outcome reach the model and inflates performance~\citep{kapoor2023leakage,varoquaux2022failures,rosenblatt2024leakage}.
The size of this inflation for grade prediction has not been measured.
\item \emph{Unknown transport to Saudi patients.} Breast cancer is the most common cancer in Saudi
women; its age-standardised incidence rose from 11.8 to 29.7 per 100{,}000 between 2001 and 2017,
when it accounted for 30.9\% of female cancers~\citep{basudan2022registry}. Patients in Arab
countries are younger and more often have hormone-receptor-negative or HER2-positive
disease than patients in Western cohorts~\citep{najjar2010arab,chouchane2013arab,elkum2007young,altamimi2010saudi}.
No grade-labelled patient-level dataset from the region is openly available, and local studies are
based on single-centre records that cannot be
shared~\citep{kussaibi2021saudi,durrani2025healthcare}. It is therefore unknown whether models built
on Western data would give correct predictions for Saudi patients~\citep{debray2015framework,steingrimsson2023transport}.
\end{enumerate}

\paragraph{Research questions}
We address these gaps with openly downloadable data only, so that every result can be reproduced by
any reader. We ask:
\textbf{RQ1}, how strongly is grade associated with each routine variable?
\textbf{RQ2}, how well do models of increasing flexibility predict grade, within the development
cohort and in independent cohorts from other countries, and are they well calibrated?
\textbf{RQ3}, how much do common analytic shortcuts inflate apparent performance?
\textbf{RQ4}, do MRI radiomic features add to the routine record?
\textbf{RQ5}, what would Western-developed models predict for patients with a Saudi case-mix?

\paragraph{Contributions}
\begin{enumerate}[leftmargin=*,itemsep=1pt,topsep=2pt]
\item A harmonised, checksummed data model of \NGradedTotal{} grade-labelled invasive breast cancers
from \NGradedCohorts{} public cohorts on three continents, built with explicit rules that exclude
in-situ disease, keep lobular and ductal carcinoma separate and allow one record per patient
(Figure~\ref{fig:framework}).
\item A model comparison that follows TRIPOD+AI~\citep{collins2024tripodai}: seven models from a
prevalence baseline to gradient boosting, nested cross-validation, multiple imputation, paired
comparisons with logistic regression, validation in five external cohorts and across datasets,
calibration with model updating, and decision curves.
\item To our knowledge, the first measurement, on the same data and models, of how much four common leakage
shortcuts inflate apparent performance for grade prediction.
\item A head-to-head test of 529 MRI radiomic features against the routine record in the same
patients.
\item A transport analysis that re-weights open cohorts to published Saudi hospital series and shows
what Western-developed models would imply for Saudi patients.
\end{enumerate}
